\chapter{Khảo sát và phân tích yêu cầu}
\label{ch:req}

\section{Khảo sát hiện trạng}

Hiện nay, phần lớn các cửa hàng thời trang trẻ em quy mô vừa và nhỏ tại Việt Nam vẫn đang
kinh doanh chủ yếu qua ba kênh: bán hàng trực tiếp tại cửa hàng, bán qua mạng xã hội
(Facebook, Instagram, TikTok) và bán qua các sàn thương mại điện tử trung gian. Qua khảo sát
thực tế, mô hình này bộc lộ nhiều hạn chế:

\begin{itemize}
  \item \textbf{Thông tin sản phẩm phân tán:} sản phẩm được đăng rải rác trên nhiều bài viết,
  không có cấu trúc danh mục rõ ràng, khó tra cứu và so sánh.
  \item \textbf{Khó quản lý tồn kho theo biến thể:} quần áo trẻ em có nhiều kích cỡ và màu sắc;
  việc theo dõi số lượng tồn của từng biến thể bằng phương pháp thủ công dễ sai sót, dẫn đến tình
  trạng bán quá số lượng tồn (oversell).
  \item \textbf{Quy trình đặt hàng thủ công:} đơn hàng được chốt qua tin nhắn, nhân viên ghi
  chép lại bằng tay, dễ nhầm lẫn và mất thời gian, khó tổng hợp doanh thu.
  \item \textbf{Phụ thuộc nền tảng trung gian:} khi bán qua sàn, cửa hàng bị động về thương hiệu,
  dữ liệu khách hàng và chính sách phí.
  \item \textbf{Thiếu công cụ chăm sóc khách hàng:} khó triển khai chương trình khuyến mãi, tích
  điểm, thu thập đánh giá một cách có hệ thống.
\end{itemize}

Từ thực trạng trên, việc xây dựng một website thương mại điện tử riêng cho thương hiệu
\brandName{} là cần thiết, nhằm tập trung hóa việc quản lý sản phẩm, tồn kho, đơn hàng và
khách hàng, đồng thời mang lại trải nghiệm mua sắm chuyên nghiệp, chủ động về thương hiệu và
dữ liệu.

\subsection{So sánh các kênh bán hàng hiện nay}

Để xác định hướng đi phù hợp cho \brandName{}, phần này so sánh các kênh bán hàng phổ biến mà
một cửa hàng thời trang trẻ em có thể lựa chọn. Mỗi kênh có những ưu, nhược điểm riêng, được
tổng hợp trong Bảng~\ref{tab:channels}.

\begin{itemize}
  \item \textbf{Bán tại cửa hàng truyền thống:} tạo được trải nghiệm trực tiếp (thử đồ, xem chất
  liệu) nhưng giới hạn về phạm vi khách hàng địa lý và thời gian mở cửa, khó mở rộng quy mô.
  \item \textbf{Bán qua mạng xã hội (Facebook, Instagram, TikTok):} dễ bắt đầu, chi phí thấp,
  tiếp cận nhanh nhưng thông tin sản phẩm phân tán theo dòng thời gian, không có cấu trúc danh
  mục, khó quản lý tồn kho và đơn hàng tập trung, phụ thuộc thuật toán hiển thị của nền tảng.
  \item \textbf{Bán qua sàn thương mại điện tử trung gian (Shopee, Lazada, TikTok Shop):} có sẵn
  lượng truy cập lớn và hạ tầng thanh toán/vận chuyển, nhưng cửa hàng phải chịu phí hoa hồng,
  cạnh tranh giá gay gắt, bị động về thương hiệu và không sở hữu dữ liệu khách hàng.
  \item \textbf{Website thương mại điện tử riêng:} chủ động hoàn toàn về thương hiệu, dữ liệu khách
  hàng, chính sách giá và khuyến mãi; quản lý sản phẩm, tồn kho, đơn hàng tập trung. Hạn chế là
  cần đầu tư xây dựng ban đầu và tự thu hút lưu lượng truy cập.
\end{itemize}

\begin{longtable}{|>{\raggedright\arraybackslash}p{0.20\linewidth}|>{\raggedright\arraybackslash}p{0.36\linewidth}|>{\raggedright\arraybackslash}p{0.34\linewidth}|}
\caption{So sánh các kênh bán hàng}\label{tab:channels}\\
\hline
\rowcolor{gray!12}\textbf{Kênh bán} & \textbf{Ưu điểm} & \textbf{Nhược điểm} \\ \hline
\endfirsthead
\hline
\rowcolor{gray!12}\textbf{Kênh bán} & \textbf{Ưu điểm} & \textbf{Nhược điểm} \\ \hline
\endhead
Cửa hàng truyền thống & Trải nghiệm trực tiếp, tạo niềm tin & Giới hạn địa lý, khó mở rộng \\ \hline
Mạng xã hội & Dễ bắt đầu, chi phí thấp, lan tỏa nhanh & Thông tin phân tán, khó quản lý kho/đơn, phụ thuộc nền tảng \\ \hline
Sàn TMĐT trung gian & Sẵn lưu lượng, hạ tầng thanh toán/vận chuyển & Phí hoa hồng, cạnh tranh giá, không sở hữu dữ liệu khách hàng \\ \hline
Website riêng & Chủ động thương hiệu, dữ liệu, khuyến mãi; quản lý tập trung & Cần đầu tư ban đầu, tự thu hút lưu lượng \\ \hline
\end{longtable}

Qua phân tích, kết hợp một website thương mại điện tử riêng với các kênh mạng xã hội là hướng đi
cân bằng: website đóng vai trò ``ngôi nhà'' chính thức của thương hiệu --- nơi quản lý tập trung
sản phẩm, đơn hàng, khuyến mãi và chăm sóc khách hàng --- trong khi mạng xã hội đảm nhiệm việc
thu hút và dẫn khách về website. Đề tài \brandName{} tập trung xây dựng thành phần website này.

\section{Tổng quan chức năng}

Hệ thống \brandName{} có hai tác nhân (actor) chính:

\begin{itemize}
  \item \textbf{Khách hàng:} người dùng cuối thực hiện mua sắm trên website.
  \item \textbf{Quản trị viên:} người quản lý vận hành toàn bộ hoạt động kinh doanh của cửa hàng
  trên hệ thống.
\end{itemize}

\subsection{Biểu đồ use case tổng quan}

Biểu đồ use case tổng quan ở Hình~\ref{fig:uc-overview} thể hiện các nhóm chức năng chính của
hai tác nhân trong hệ thống.

\fig{figures/diagrams/uc-overview.pdf}{0.62}{\caption{Biểu đồ use case tổng quan}\label{fig:uc-overview}}

\subsection{Biểu đồ các ca sử dụng khách hàng}

Hình~\ref{fig:uc-customer} mô tả chi tiết các ca sử dụng của tác nhân Khách hàng.

\fig{figures/diagrams/uc-customer.pdf}{0.95}{\caption{Biểu đồ use case khách hàng}\label{fig:uc-customer}}

\subsection{Biểu đồ các ca sử dụng quản trị viên}

Hình~\ref{fig:uc-admin} mô tả chi tiết các ca sử dụng của tác nhân Quản trị viên.

\fig{figures/diagrams/uc-admin.pdf}{0.9}{\caption{Biểu đồ use case quản trị viên}\label{fig:uc-admin}}

\subsection{Phân tích quy trình nghiệp vụ}

Phần này phân tích ba quy trình nghiệp vụ tiêu biểu của hệ thống bằng biểu đồ hoạt động.

\textbf{Quy trình đặt hàng và thanh toán.} Đây là quy trình cốt lõi của hệ thống, bao gồm việc
kiểm tra tồn kho, tạo đơn hàng và xử lý thanh toán theo hai hình thức COD và PayOS
(Hình~\ref{fig:bp-order}).

\fig{figures/diagrams/bp-order.pdf}{0.62}{\caption{Quy trình nghiệp vụ Đặt hàng và thanh toán}\label{fig:bp-order}}

\textbf{Quy trình quản lý sản phẩm.} Mô tả các bước quản trị viên thêm mới hoặc cập nhật sản
phẩm cùng các biến thể và tồn kho (Hình~\ref{fig:bp-product}).

\fig{figures/diagrams/bp-product.pdf}{0.55}{\caption{Quy trình nghiệp vụ Quản lý sản phẩm}\label{fig:bp-product}}

\textbf{Quy trình xử lý đổi/trả hàng.} Mô tả luồng khách hàng gửi yêu cầu đổi/trả và quản trị
viên xử lý, bao gồm cả nghiệp vụ hoàn tiền và hoàn trả tồn kho (Hình~\ref{fig:bp-return}).

\fig{figures/diagrams/bp-return.pdf}{0.6}{\caption{Quy trình nghiệp vụ Xử lý đổi/trả hàng}\label{fig:bp-return}}

\section{Đặc tả chức năng}

\subsection{Các ca sử dụng của người mua}

Bảng~\ref{tab:uc-customer-list} liệt kê danh sách các ca sử dụng của tác nhân Khách hàng.

\begin{longtable}{|c|>{\raggedright\arraybackslash}p{0.30\linewidth}|>{\raggedright\arraybackslash}p{0.50\linewidth}|}
\caption{Danh sách use case của Khách hàng}\label{tab:uc-customer-list}\\
\hline
\rowcolor{gray!12}\textbf{Mã} & \textbf{Tên use case} & \textbf{Mô tả ngắn} \\ \hline
\endfirsthead
\hline
\rowcolor{gray!12}\textbf{Mã} & \textbf{Tên use case} & \textbf{Mô tả ngắn} \\ \hline
\endhead
UC01 & Đăng ký & Tạo tài khoản khách hàng mới \\ \hline
UC02 & Đăng nhập & Đăng nhập vào hệ thống \\ \hline
UC03 & Đăng xuất & Thoát khỏi phiên đăng nhập \\ \hline
UC04 & Xem danh sách sản phẩm & Duyệt sản phẩm theo danh mục \\ \hline
UC05 & Tìm kiếm sản phẩm & Tìm sản phẩm theo từ khóa \\ \hline
UC06 & Lọc sản phẩm & Lọc theo độ tuổi, giới tính, giá, kích cỡ, màu \\ \hline
UC07 & Xem chi tiết sản phẩm & Xem thông tin, biến thể, đánh giá \\ \hline
UC08 & Quản lý giỏ hàng & Thêm/sửa/xóa sản phẩm trong giỏ \\ \hline
UC09 & Quản lý danh sách yêu thích & Lưu và bỏ lưu sản phẩm yêu thích \\ \hline
UC10 & So sánh sản phẩm & So sánh thông số nhiều sản phẩm \\ \hline
UC11 & Tạo đơn hàng & Khởi tạo đơn từ giỏ hàng \\ \hline
UC12 & Áp dụng mã giảm giá & Nhập và áp dụng voucher \\ \hline
UC13 & Đặt hàng & Xác nhận đặt đơn hàng \\ \hline
UC14 & Thanh toán đơn hàng & Thanh toán COD hoặc PayOS \\ \hline
UC15 & Hủy đơn hàng & Hủy đơn khi còn cho phép \\ \hline
UC16 & Xem lịch sử đơn hàng & Theo dõi đơn và trạng thái \\ \hline
UC17 & Yêu cầu đổi/trả hàng & Gửi yêu cầu đổi/trả đơn đã hoàn thành \\ \hline
UC18 & Đánh giá, bình luận sản phẩm & Gửi đánh giá và bình luận \\ \hline
UC19 & Quản lý thông tin cá nhân & Xem và cập nhật hồ sơ \\ \hline
UC20 & Quản lý địa chỉ giao hàng & Thêm/sửa/xóa địa chỉ \\ \hline
\end{longtable}

\subsection{Các ca sử dụng của người quản trị}

Bảng~\ref{tab:uc-admin-list} liệt kê danh sách các ca sử dụng của tác nhân Quản trị viên.

\begin{longtable}{|c|>{\raggedright\arraybackslash}p{0.30\linewidth}|>{\raggedright\arraybackslash}p{0.50\linewidth}|}
\caption{Danh sách use case của Quản trị viên}\label{tab:uc-admin-list}\\
\hline
\rowcolor{gray!12}\textbf{Mã} & \textbf{Tên use case} & \textbf{Mô tả ngắn} \\ \hline
\endfirsthead
\hline
\rowcolor{gray!12}\textbf{Mã} & \textbf{Tên use case} & \textbf{Mô tả ngắn} \\ \hline
\endhead
UC21 & Đăng nhập quản trị & Đăng nhập với quyền quản trị \\ \hline
UC22 & Quản lý danh mục sản phẩm & Thêm/sửa/xóa danh mục \\ \hline
UC23 & Quản lý sản phẩm \& biến thể & Quản lý sản phẩm, biến thể, hình ảnh \\ \hline
UC24 & Quản lý tồn kho & Điều chỉnh và theo dõi tồn kho \\ \hline
UC25 & Quản lý đơn hàng & Xác nhận, cập nhật trạng thái, mã vận đơn \\ \hline
UC26 & Xử lý đổi/trả hàng & Duyệt, từ chối, hoàn tiền \\ \hline
UC27 & Quản lý mã giảm giá & Thêm/sửa/xóa voucher \\ \hline
UC28 & Quản lý flash sale & Tạo và quản lý chương trình flash sale \\ \hline
UC29 & Quản lý banner \& nội dung & Quản lý banner trang chủ \\ \hline
UC30 & Quản lý tài khoản khách hàng & Xem danh sách và chi tiết khách hàng \\ \hline
UC31 & Kiểm duyệt đánh giá & Duyệt/từ chối/phản hồi đánh giá \\ \hline
UC32 & Thống kê, báo cáo doanh thu & Xem báo cáo và xuất Excel \\ \hline
\end{longtable}

\section{Đặc tả chi tiết các ca sử dụng của Khách hàng}

\subsection{Đặc tả use case ``Đăng ký''}

\begin{ucspec}{Đặc tả use case ``Đăng ký''}
\textbf{Mã use case} & UC01 \\ \hline
\textbf{Tên use case} & Đăng ký \\ \hline
\textbf{Tác nhân} & Khách truy cập \\ \hline
\textbf{Mục đích} & Cho phép người dùng tạo tài khoản mới trên hệ thống \\ \hline
\textbf{Tiền điều kiện} & Người dùng chưa có tài khoản \\ \hline
\textbf{Hậu điều kiện} & Tài khoản mới được tạo; người dùng có thể đăng nhập \\ \hline
\textbf{Luồng sự kiện chính} &
1. Khách chọn chức năng ``Đăng ký''. \newline
2. Hệ thống hiển thị biểu mẫu đăng ký. \newline
3. Khách nhập email, mật khẩu, họ tên và gửi biểu mẫu. \newline
4. Hệ thống kiểm tra hợp lệ và email chưa tồn tại. \newline
5. Hệ thống tạo tài khoản qua Supabase Auth và gửi email xác nhận. \newline
6. Khách xác nhận email và được chuyển tới trang đăng nhập. \\ \hline
\textbf{Luồng thay thế} &
4a. Dữ liệu không hợp lệ: hệ thống báo lỗi và yêu cầu nhập lại. \newline
4b. Email đã tồn tại: hệ thống thông báo và đề nghị đăng nhập. \\ \hline
\end{ucspec}

Dữ liệu đầu vào của use case ``Đăng ký'' được mô tả trong Bảng~\ref{tab:reg-input}.

\begin{longtable}{|c|>{\raggedright\arraybackslash}p{0.20\linewidth}|c|>{\raggedright\arraybackslash}p{0.34\linewidth}|}
\caption{Dữ liệu đầu vào use case ``Đăng ký''}\label{tab:reg-input}\\
\hline
\rowcolor{gray!12}\textbf{STT} & \textbf{Trường} & \textbf{Bắt buộc} & \textbf{Điều kiện hợp lệ} \\ \hline
\endfirsthead
\hline
\rowcolor{gray!12}\textbf{STT} & \textbf{Trường} & \textbf{Bắt buộc} & \textbf{Điều kiện hợp lệ} \\ \hline
\endhead
1 & Họ tên & Có & Không để trống \\ \hline
2 & Email & Có & Đúng định dạng email, chưa tồn tại \\ \hline
3 & Mật khẩu & Có & Tối thiểu 6 ký tự \\ \hline
4 & Xác nhận mật khẩu & Có & Trùng khớp với mật khẩu \\ \hline
\end{longtable}

\subsection{Đặc tả use case ``Đăng nhập''}

\begin{ucspec}{Đặc tả use case ``Đăng nhập''}
\textbf{Mã use case} & UC02 \\ \hline
\textbf{Tên use case} & Đăng nhập \\ \hline
\textbf{Tác nhân} & Khách hàng \\ \hline
\textbf{Mục đích} & Xác thực và bắt đầu phiên làm việc của người dùng \\ \hline
\textbf{Tiền điều kiện} & Người dùng đã có tài khoản \\ \hline
\textbf{Hậu điều kiện} & Người dùng đăng nhập thành công, phiên được thiết lập \\ \hline
\textbf{Luồng sự kiện chính} &
1. Khách chọn ``Đăng nhập''. \newline
2. Hệ thống hiển thị biểu mẫu đăng nhập. \newline
3. Khách nhập email và mật khẩu. \newline
4. Hệ thống xác thực thông tin qua Supabase Auth. \newline
5. Đăng nhập thành công, chuyển về trang trước đó hoặc trang chủ. \\ \hline
\textbf{Luồng thay thế} &
4a. Sai email hoặc mật khẩu: hệ thống thông báo lỗi. \newline
4b. Vượt quá số lần thử cho phép: hệ thống tạm giới hạn truy cập (rate limit). \\ \hline
\end{ucspec}

\subsection{Đặc tả use case ``Tìm kiếm sản phẩm''}

\begin{ucspec}{Đặc tả use case ``Tìm kiếm sản phẩm''}
\textbf{Mã use case} & UC05 \\ \hline
\textbf{Tên use case} & Tìm kiếm sản phẩm \\ \hline
\textbf{Tác nhân} & Khách hàng \\ \hline
\textbf{Mục đích} & Tìm sản phẩm theo từ khóa \\ \hline
\textbf{Tiền điều kiện} & Không \\ \hline
\textbf{Hậu điều kiện} & Hiển thị danh sách sản phẩm phù hợp \\ \hline
\textbf{Luồng sự kiện chính} &
1. Khách nhập từ khóa vào ô tìm kiếm. \newline
2. Hệ thống chuẩn hóa từ khóa, loại bỏ dấu (\texttt{unaccent}). \newline
3. Hệ thống truy vấn tìm kiếm toàn văn trên cơ sở dữ liệu. \newline
4. Hệ thống hiển thị danh sách sản phẩm phù hợp, sắp theo độ liên quan. \\ \hline
\textbf{Luồng thay thế} &
4a. Không có kết quả: hệ thống hiển thị thông báo ``không tìm thấy sản phẩm''. \\ \hline
\end{ucspec}

\subsection{Đặc tả use case ``Lọc danh sách sản phẩm''}

\begin{ucspec}{Đặc tả use case ``Lọc danh sách sản phẩm''}
\textbf{Mã use case} & UC06 \\ \hline
\textbf{Tên use case} & Lọc danh sách sản phẩm \\ \hline
\textbf{Tác nhân} & Khách hàng \\ \hline
\textbf{Mục đích} & Thu hẹp danh sách theo nhiều tiêu chí \\ \hline
\textbf{Tiền điều kiện} & Đang ở trang danh sách sản phẩm \\ \hline
\textbf{Hậu điều kiện} & Danh sách được lọc theo tiêu chí đã chọn \\ \hline
\textbf{Luồng sự kiện chính} &
1. Khách chọn các tiêu chí lọc: danh mục, độ tuổi, giới tính, khoảng giá, kích cỡ, màu sắc. \newline
2. Hệ thống cập nhật tham số lọc trên URL. \newline
3. Hệ thống truy vấn và hiển thị danh sách sản phẩm thỏa điều kiện. \\ \hline
\textbf{Luồng thay thế} &
3a. Không có sản phẩm phù hợp: hiển thị thông báo và gợi ý bỏ bớt bộ lọc. \\ \hline
\end{ucspec}

\subsection{Đặc tả use case ``Quản lý giỏ hàng''}

\begin{ucspec}{Đặc tả use case ``Quản lý giỏ hàng''}
\textbf{Mã use case} & UC08 \\ \hline
\textbf{Tên use case} & Quản lý giỏ hàng \\ \hline
\textbf{Tác nhân} & Khách hàng \\ \hline
\textbf{Mục đích} & Thêm, cập nhật, xóa sản phẩm trong giỏ trước khi đặt hàng \\ \hline
\textbf{Tiền điều kiện} & Không \\ \hline
\textbf{Hậu điều kiện} & Giỏ hàng được cập nhật \\ \hline
\textbf{Luồng sự kiện chính} &
1. Khách chọn biến thể (kích cỡ, màu) và thêm sản phẩm vào giỏ. \newline
2. Hệ thống lưu giỏ hàng (Zustand, lưu cục bộ trên trình duyệt). \newline
3. Khách có thể tăng/giảm số lượng hoặc xóa sản phẩm khỏi giỏ. \newline
4. Hệ thống tính lại tổng tiền tạm tính. \\ \hline
\textbf{Luồng thay thế} &
3a. Số lượng vượt quá tồn kho: hệ thống giới hạn theo số lượng còn lại. \\ \hline
\end{ucspec}

\subsection{Đặc tả use case ``Quản lý danh sách yêu thích''}

\begin{ucspec}{Đặc tả use case ``Quản lý danh sách yêu thích''}
\textbf{Mã use case} & UC09 \\ \hline
\textbf{Tên use case} & Quản lý danh sách yêu thích \\ \hline
\textbf{Tác nhân} & Khách hàng (đã đăng nhập) \\ \hline
\textbf{Mục đích} & Lưu lại các sản phẩm quan tâm để xem/mua sau \\ \hline
\textbf{Tiền điều kiện} & Khách đã đăng nhập \\ \hline
\textbf{Hậu điều kiện} & Danh sách yêu thích được cập nhật \\ \hline
\textbf{Luồng sự kiện chính} &
1. Khách nhấn biểu tượng trái tim trên sản phẩm. \newline
2. Hệ thống lưu sản phẩm vào danh sách yêu thích của tài khoản. \newline
3. Khách có thể xem trang yêu thích và bỏ lưu sản phẩm. \\ \hline
\textbf{Luồng thay thế} &
1a. Khách chưa đăng nhập: hệ thống yêu cầu đăng nhập. \\ \hline
\end{ucspec}

\subsection{Đặc tả use case ``So sánh sản phẩm''}

\begin{ucspec}{Đặc tả use case ``So sánh sản phẩm''}
\textbf{Mã use case} & UC10 \\ \hline
\textbf{Tên use case} & So sánh sản phẩm \\ \hline
\textbf{Tác nhân} & Khách hàng \\ \hline
\textbf{Mục đích} & So sánh thông số của nhiều sản phẩm cạnh nhau \\ \hline
\textbf{Tiền điều kiện} & Không \\ \hline
\textbf{Hậu điều kiện} & Hiển thị bảng so sánh các sản phẩm đã chọn \\ \hline
\textbf{Luồng sự kiện chính} &
1. Khách thêm sản phẩm vào danh sách so sánh (tối đa 4 sản phẩm). \newline
2. Hệ thống hiển thị thanh so sánh nổi. \newline
3. Khách mở trang so sánh; hệ thống hiển thị bảng thông số đối chiếu. \\ \hline
\textbf{Luồng thay thế} &
1a. Vượt quá 4 sản phẩm: hệ thống thông báo giới hạn. \\ \hline
\end{ucspec}

\subsection{Đặc tả use case ``Tạo đơn hàng''}

\begin{ucspec}{Đặc tả use case ``Tạo đơn hàng''}
\textbf{Mã use case} & UC11 \\ \hline
\textbf{Tên use case} & Tạo đơn hàng \\ \hline
\textbf{Tác nhân} & Khách hàng \\ \hline
\textbf{Mục đích} & Khởi tạo đơn hàng từ các sản phẩm trong giỏ \\ \hline
\textbf{Tiền điều kiện} & Giỏ hàng có ít nhất một sản phẩm; khách đã đăng nhập \\ \hline
\textbf{Hậu điều kiện} & Đơn hàng được tạo ở trạng thái chờ xác nhận \\ \hline
\textbf{Luồng sự kiện chính} &
1. Khách vào trang thanh toán. \newline
2. Khách chọn địa chỉ giao hàng và phương thức thanh toán. \newline
3. Hệ thống tính phí vận chuyển và tổng tiền. \newline
4. Hệ thống kiểm tra tồn kho và tạo đơn hàng trong một giao dịch. \\ \hline
\textbf{Luồng thay thế} &
4a. Một sản phẩm hết hàng: hệ thống hủy giao dịch và thông báo. \\ \hline
\end{ucspec}

\subsection{Đặc tả use case ``Đặt hàng''}

\begin{ucspec}{Đặc tả use case ``Đặt hàng''}
\textbf{Mã use case} & UC13 \\ \hline
\textbf{Tên use case} & Đặt hàng \\ \hline
\textbf{Tác nhân} & Khách hàng \\ \hline
\textbf{Mục đích} & Xác nhận và hoàn tất việc đặt đơn hàng \\ \hline
\textbf{Tiền điều kiện} & Đơn hàng đã được tạo và hợp lệ \\ \hline
\textbf{Hậu điều kiện} & Đơn hàng được ghi nhận; tồn kho được trừ tương ứng \\ \hline
\textbf{Luồng sự kiện chính} &
1. Khách xác nhận đặt hàng. \newline
2. Hệ thống lưu đơn hàng, trừ tồn kho và sinh mã đơn hàng. \newline
3. Nếu chọn PayOS, hệ thống tạo liên kết thanh toán. \newline
4. Hệ thống gửi email xác nhận đặt hàng và hiển thị trang xác nhận. \\ \hline
\textbf{Luồng thay thế} &
3a. Tạo liên kết PayOS thất bại: hệ thống hủy đơn, hoàn tồn kho và báo lỗi. \\ \hline
\end{ucspec}

\subsection{Đặc tả use case ``Thanh toán đơn hàng''}

\begin{ucspec}{Đặc tả use case ``Thanh toán đơn hàng''}
\textbf{Mã use case} & UC14 \\ \hline
\textbf{Tên use case} & Thanh toán đơn hàng \\ \hline
\textbf{Tác nhân} & Khách hàng, Cổng PayOS \\ \hline
\textbf{Mục đích} & Thanh toán trực tuyến cho đơn hàng \\ \hline
\textbf{Tiền điều kiện} & Đơn hàng PayOS đã được tạo và chưa thanh toán \\ \hline
\textbf{Hậu điều kiện} & Đơn hàng được đánh dấu đã thanh toán \\ \hline
\textbf{Luồng sự kiện chính} &
1. Khách được chuyển tới trang thanh toán của PayOS. \newline
2. Khách quét mã QR / chuyển khoản để thanh toán. \newline
3. PayOS gọi webhook thông báo kết quả; hệ thống xác thực chữ ký HMAC. \newline
4. Hệ thống cập nhật trạng thái ``đã thanh toán'' và gửi email xác nhận. \\ \hline
\textbf{Luồng thay thế} &
4a. Quá thời hạn (24 giờ) chưa thanh toán: tác vụ định kỳ hủy đơn và hoàn tồn kho. \\ \hline
\end{ucspec}

\subsection{Đặc tả use case ``Hủy đơn hàng''}

\begin{ucspec}{Đặc tả use case ``Hủy đơn hàng''}
\textbf{Mã use case} & UC15 \\ \hline
\textbf{Tên use case} & Hủy đơn hàng \\ \hline
\textbf{Tác nhân} & Khách hàng \\ \hline
\textbf{Mục đích} & Hủy đơn hàng khi còn được phép \\ \hline
\textbf{Tiền điều kiện} & Đơn ở trạng thái chờ xác nhận hoặc chưa thanh toán \\ \hline
\textbf{Hậu điều kiện} & Đơn chuyển sang trạng thái đã hủy; tồn kho được hoàn trả \\ \hline
\textbf{Luồng sự kiện chính} &
1. Khách chọn đơn hàng và nhấn ``Hủy đơn''. \newline
2. Hệ thống kiểm tra điều kiện cho phép hủy. \newline
3. Hệ thống cập nhật trạng thái và hoàn trả tồn kho. \\ \hline
\textbf{Luồng thay thế} &
2a. Đơn không cho phép hủy: hệ thống thông báo từ chối. \\ \hline
\end{ucspec}

\subsection{Đặc tả use case ``Xem lịch sử đơn hàng''}

\begin{ucspec}{Đặc tả use case ``Xem lịch sử đơn hàng''}
\textbf{Mã use case} & UC16 \\ \hline
\textbf{Tên use case} & Xem lịch sử đơn hàng \\ \hline
\textbf{Tác nhân} & Khách hàng \\ \hline
\textbf{Mục đích} & Theo dõi các đơn hàng đã đặt và trạng thái của chúng \\ \hline
\textbf{Tiền điều kiện} & Khách đã đăng nhập \\ \hline
\textbf{Hậu điều kiện} & Hiển thị danh sách và chi tiết đơn hàng \\ \hline
\textbf{Luồng sự kiện chính} &
1. Khách vào mục ``Đơn hàng của tôi''. \newline
2. Hệ thống hiển thị danh sách đơn theo thứ tự thời gian. \newline
3. Khách chọn một đơn để xem chi tiết và trạng thái giao hàng. \\ \hline
\textbf{Luồng thay thế} & --- \\ \hline
\end{ucspec}

\subsection{Đặc tả use case ``Yêu cầu đổi/trả hàng''}

\begin{ucspec}{Đặc tả use case ``Yêu cầu đổi/trả hàng''}
\textbf{Mã use case} & UC17 \\ \hline
\textbf{Tên use case} & Yêu cầu đổi/trả hàng \\ \hline
\textbf{Tác nhân} & Khách hàng \\ \hline
\textbf{Mục đích} & Gửi yêu cầu đổi hoặc trả hàng cho đơn đã hoàn thành \\ \hline
\textbf{Tiền điều kiện} & Đơn hàng ở trạng thái đã hoàn thành \\ \hline
\textbf{Hậu điều kiện} & Yêu cầu đổi/trả được tạo, chờ quản trị xử lý \\ \hline
\textbf{Luồng sự kiện chính} &
1. Khách mở đơn đã hoàn thành và chọn ``Yêu cầu đổi/trả''. \newline
2. Khách nhập lý do và gửi yêu cầu. \newline
3. Hệ thống tạo yêu cầu ở trạng thái ``đã yêu cầu''. \\ \hline
\textbf{Luồng thay thế} &
1a. Đơn chưa hoàn thành: hệ thống không cho phép tạo yêu cầu. \\ \hline
\end{ucspec}

\subsection{Đặc tả use case ``Đánh giá, bình luận sản phẩm''}

\begin{ucspec}{Đặc tả use case ``Đánh giá, bình luận sản phẩm''}
\textbf{Mã use case} & UC18 \\ \hline
\textbf{Tên use case} & Đánh giá, bình luận sản phẩm \\ \hline
\textbf{Tác nhân} & Khách hàng \\ \hline
\textbf{Mục đích} & Gửi đánh giá sao, nhận xét và hình ảnh về sản phẩm \\ \hline
\textbf{Tiền điều kiện} & Khách đã đăng nhập \\ \hline
\textbf{Hậu điều kiện} & Đánh giá được lưu ở trạng thái chờ duyệt \\ \hline
\textbf{Luồng sự kiện chính} &
1. Khách mở chi tiết sản phẩm, chọn tab đánh giá. \newline
2. Khách chấm số sao, nhập nhận xét và tải tối đa 5 ảnh. \newline
3. Hệ thống lưu đánh giá ở trạng thái ``chờ duyệt''. \newline
4. Sau khi được quản trị duyệt, đánh giá hiển thị công khai và cập nhật điểm trung bình. \\ \hline
\textbf{Luồng thay thế} &
2a. Khách đã đánh giá sản phẩm trước đó: hệ thống cho phép cập nhật đánh giá cũ. \\ \hline
\end{ucspec}

Dữ liệu đầu vào của use case ``Đăng nhập'' được mô tả trong Bảng~\ref{tab:login-input}.

\begin{longtable}{|c|>{\raggedright\arraybackslash}p{0.20\linewidth}|c|>{\raggedright\arraybackslash}p{0.34\linewidth}|}
\caption{Dữ liệu đầu vào use case ``Đăng nhập''}\label{tab:login-input}\\
\hline
\rowcolor{gray!12}\textbf{STT} & \textbf{Trường} & \textbf{Bắt buộc} & \textbf{Điều kiện hợp lệ} \\ \hline
\endfirsthead
\hline
\rowcolor{gray!12}\textbf{STT} & \textbf{Trường} & \textbf{Bắt buộc} & \textbf{Điều kiện hợp lệ} \\ \hline
\endhead
1 & Email & Có & Đúng định dạng email \\ \hline
2 & Mật khẩu & Có & Không để trống \\ \hline
\end{longtable}

\subsection{Đặc tả use case ``Đăng xuất''}

\begin{ucspec}{Đặc tả use case ``Đăng xuất''}
\textbf{Mã use case} & UC03 \\ \hline
\textbf{Tên use case} & Đăng xuất \\ \hline
\textbf{Tác nhân} & Khách hàng \\ \hline
\textbf{Mục đích} & Kết thúc phiên làm việc của người dùng \\ \hline
\textbf{Tiền điều kiện} & Người dùng đang đăng nhập \\ \hline
\textbf{Hậu điều kiện} & Phiên đăng nhập được kết thúc \\ \hline
\textbf{Luồng sự kiện chính} &
1. Khách chọn chức năng ``Đăng xuất''. \newline
2. Hệ thống xóa phiên đăng nhập và chuyển về trang chủ. \\ \hline
\textbf{Luồng thay thế} & --- \\ \hline
\end{ucspec}

\subsection{Đặc tả use case ``Xem danh sách sản phẩm''}

\begin{ucspec}{Đặc tả use case ``Xem danh sách sản phẩm''}
\textbf{Mã use case} & UC04 \\ \hline
\textbf{Tên use case} & Xem danh sách sản phẩm \\ \hline
\textbf{Tác nhân} & Khách hàng \\ \hline
\textbf{Mục đích} & Duyệt danh sách sản phẩm theo danh mục \\ \hline
\textbf{Tiền điều kiện} & Không \\ \hline
\textbf{Hậu điều kiện} & Hiển thị danh sách sản phẩm \\ \hline
\textbf{Luồng sự kiện chính} &
1. Khách truy cập trang sản phẩm hoặc chọn một danh mục. \newline
2. Hệ thống truy vấn các sản phẩm đang bán và phân trang. \newline
3. Hệ thống hiển thị danh sách kèm hình ảnh, tên, giá và nhãn (mới/sale). \\ \hline
\textbf{Luồng thay thế} &
2a. Danh mục không có sản phẩm: hiển thị thông báo phù hợp. \\ \hline
\end{ucspec}

\subsection{Đặc tả use case ``Xem chi tiết sản phẩm''}

\begin{ucspec}{Đặc tả use case ``Xem chi tiết sản phẩm''}
\textbf{Mã use case} & UC07 \\ \hline
\textbf{Tên use case} & Xem chi tiết sản phẩm \\ \hline
\textbf{Tác nhân} & Khách hàng \\ \hline
\textbf{Mục đích} & Xem thông tin đầy đủ của một sản phẩm \\ \hline
\textbf{Tiền điều kiện} & Không \\ \hline
\textbf{Hậu điều kiện} & Hiển thị chi tiết sản phẩm \\ \hline
\textbf{Luồng sự kiện chính} &
1. Khách chọn một sản phẩm từ danh sách. \newline
2. Hệ thống hiển thị hình ảnh, mô tả, giá, các biến thể (kích cỡ, màu), bảng quy đổi size,
thông số và danh sách đánh giá. \newline
3. Hệ thống gợi ý các sản phẩm tương tự. \\ \hline
\textbf{Luồng thay thế} &
2a. Sản phẩm đang trong flash sale: hiển thị giá khuyến mãi và đồng hồ đếm ngược. \\ \hline
\end{ucspec}

\subsection{Đặc tả use case ``Áp dụng mã giảm giá''}

\begin{ucspec}{Đặc tả use case ``Áp dụng mã giảm giá''}
\textbf{Mã use case} & UC12 \\ \hline
\textbf{Tên use case} & Áp dụng mã giảm giá \\ \hline
\textbf{Tác nhân} & Khách hàng \\ \hline
\textbf{Mục đích} & Áp dụng mã giảm giá cho đơn hàng \\ \hline
\textbf{Tiền điều kiện} & Đang ở trang thanh toán, giỏ có sản phẩm \\ \hline
\textbf{Hậu điều kiện} & Mã hợp lệ được áp dụng, tổng tiền được giảm tương ứng \\ \hline
\textbf{Luồng sự kiện chính} &
1. Khách nhập mã giảm giá và nhấn ``Áp dụng''. \newline
2. Hệ thống kiểm tra mã: tồn tại, đang hiệu lực, đạt giá trị đơn tối thiểu, còn lượt sử dụng. \newline
3. Hệ thống tính mức giảm và cập nhật tổng tiền tạm tính. \\ \hline
\textbf{Luồng thay thế} &
2a. Mã không hợp lệ/hết hạn/hết lượt: hệ thống thông báo lỗi tương ứng. \\ \hline
\end{ucspec}

\subsection{Đặc tả use case ``Quản lý thông tin cá nhân''}

\begin{ucspec}{Đặc tả use case ``Quản lý thông tin cá nhân''}
\textbf{Mã use case} & UC19 \\ \hline
\textbf{Tên use case} & Quản lý thông tin cá nhân \\ \hline
\textbf{Tác nhân} & Khách hàng \\ \hline
\textbf{Mục đích} & Xem và cập nhật hồ sơ cá nhân \\ \hline
\textbf{Tiền điều kiện} & Khách đã đăng nhập \\ \hline
\textbf{Hậu điều kiện} & Thông tin hồ sơ được cập nhật \\ \hline
\textbf{Luồng sự kiện chính} &
1. Khách vào trang tài khoản. \newline
2. Khách cập nhật họ tên, số điện thoại. \newline
3. Hệ thống kiểm tra hợp lệ và lưu thay đổi. \\ \hline
\textbf{Luồng thay thế} &
3a. Dữ liệu không hợp lệ (số điện thoại sai định dạng): hệ thống báo lỗi. \\ \hline
\end{ucspec}

Dữ liệu đầu vào của use case ``Quản lý thông tin cá nhân'' được mô tả trong
Bảng~\ref{tab:profile-input}.

\begin{longtable}{|c|>{\raggedright\arraybackslash}p{0.20\linewidth}|c|>{\raggedright\arraybackslash}p{0.34\linewidth}|}
\caption{Dữ liệu đầu vào use case ``Quản lý thông tin cá nhân''}\label{tab:profile-input}\\
\hline
\rowcolor{gray!12}\textbf{STT} & \textbf{Trường} & \textbf{Bắt buộc} & \textbf{Điều kiện hợp lệ} \\ \hline
\endfirsthead
\hline
\rowcolor{gray!12}\textbf{STT} & \textbf{Trường} & \textbf{Bắt buộc} & \textbf{Điều kiện hợp lệ} \\ \hline
\endhead
1 & Họ tên & Có & Từ 2 đến 100 ký tự \\ \hline
2 & Số điện thoại & Không & Định dạng \texttt{0} hoặc \texttt{+84} + 9 chữ số \\ \hline
\end{longtable}

\subsection{Đặc tả use case ``Quản lý địa chỉ giao hàng''}

\begin{ucspec}{Đặc tả use case ``Quản lý địa chỉ giao hàng''}
\textbf{Mã use case} & UC20 \\ \hline
\textbf{Tên use case} & Quản lý địa chỉ giao hàng \\ \hline
\textbf{Tác nhân} & Khách hàng \\ \hline
\textbf{Mục đích} & Thêm, sửa, xóa và đặt mặc định địa chỉ giao hàng \\ \hline
\textbf{Tiền điều kiện} & Khách đã đăng nhập \\ \hline
\textbf{Hậu điều kiện} & Sổ địa chỉ được cập nhật \\ \hline
\textbf{Luồng sự kiện chính} &
1. Khách vào mục địa chỉ trong trang tài khoản. \newline
2. Khách thêm/sửa địa chỉ: nhập tên người nhận, số điện thoại, tỉnh/quận/phường, địa chỉ chi
tiết; có thể đặt làm mặc định. \newline
3. Hệ thống kiểm tra hợp lệ và lưu địa chỉ. \\ \hline
\textbf{Luồng thay thế} &
3a. Dữ liệu không hợp lệ: hệ thống báo lỗi và yêu cầu sửa. \\ \hline
\end{ucspec}

Dữ liệu đầu vào của use case ``Quản lý địa chỉ giao hàng'' được mô tả trong
Bảng~\ref{tab:address-input}.

\begin{longtable}{|c|>{\raggedright\arraybackslash}p{0.24\linewidth}|c|>{\raggedright\arraybackslash}p{0.30\linewidth}|}
\caption{Dữ liệu đầu vào use case ``Quản lý địa chỉ giao hàng''}\label{tab:address-input}\\
\hline
\rowcolor{gray!12}\textbf{STT} & \textbf{Trường} & \textbf{Bắt buộc} & \textbf{Điều kiện hợp lệ} \\ \hline
\endfirsthead
\hline
\rowcolor{gray!12}\textbf{STT} & \textbf{Trường} & \textbf{Bắt buộc} & \textbf{Điều kiện hợp lệ} \\ \hline
\endhead
1 & Tên người nhận & Có & Từ 2 đến 100 ký tự \\ \hline
2 & Số điện thoại & Có & Định dạng \texttt{0}/\texttt{+84} + 9 chữ số \\ \hline
3 & Tỉnh/thành phố & Có & Chọn từ danh sách \\ \hline
4 & Quận/huyện & Có & Chọn từ danh sách \\ \hline
5 & Phường/xã & Có & Chọn từ danh sách \\ \hline
6 & Địa chỉ chi tiết & Có & Từ 5 đến 255 ký tự \\ \hline
7 & Đặt làm mặc định & Không & Đúng/Sai \\ \hline
\end{longtable}

\section{Đặc tả chi tiết các ca sử dụng của Quản trị viên}

\subsection{Đặc tả use case ``Quản lý sản phẩm''}

\begin{ucspec}{Đặc tả use case ``Quản lý sản phẩm''}
\textbf{Mã use case} & UC23 \\ \hline
\textbf{Tên use case} & Quản lý sản phẩm \& biến thể \\ \hline
\textbf{Tác nhân} & Quản trị viên \\ \hline
\textbf{Mục đích} & Thêm, sửa, xóa sản phẩm cùng biến thể và hình ảnh \\ \hline
\textbf{Tiền điều kiện} & Quản trị viên đã đăng nhập \\ \hline
\textbf{Hậu điều kiện} & Dữ liệu sản phẩm được cập nhật \\ \hline
\textbf{Luồng sự kiện chính} &
1. Quản trị viên vào trang quản lý sản phẩm. \newline
2. Chọn thêm mới hoặc chỉnh sửa một sản phẩm. \newline
3. Nhập thông tin, tải hình ảnh và khai báo biến thể (kích cỡ, màu, SKU, tồn kho). \newline
4. Hệ thống kiểm tra hợp lệ và lưu vào cơ sở dữ liệu. \\ \hline
\textbf{Luồng thay thế} &
4a. Dữ liệu không hợp lệ: hệ thống báo lỗi và yêu cầu sửa lại. \\ \hline
\end{ucspec}

\subsection{Đặc tả use case ``Quản lý đơn hàng''}

\begin{ucspec}{Đặc tả use case ``Quản lý đơn hàng''}
\textbf{Mã use case} & UC25 \\ \hline
\textbf{Tên use case} & Quản lý đơn hàng \\ \hline
\textbf{Tác nhân} & Quản trị viên \\ \hline
\textbf{Mục đích} & Xác nhận và cập nhật trạng thái xử lý đơn hàng \\ \hline
\textbf{Tiền điều kiện} & Quản trị viên đã đăng nhập \\ \hline
\textbf{Hậu điều kiện} & Trạng thái đơn hàng được cập nhật \\ \hline
\textbf{Luồng sự kiện chính} &
1. Quản trị viên xem danh sách đơn hàng. \newline
2. Chọn một đơn để xem chi tiết. \newline
3. Cập nhật trạng thái (xác nhận, đóng gói, giao hàng, hoàn thành) và nhập mã vận đơn. \newline
4. Hệ thống lưu thay đổi và gửi email thông báo cho khách. \newline
5. Quản trị viên có thể in phiếu giao hàng. \\ \hline
\textbf{Luồng thay thế} &
3a. Hủy đơn: hệ thống hoàn trả tồn kho cho đơn bị hủy. \\ \hline
\end{ucspec}

\subsection{Đặc tả use case ``Quản lý mã giảm giá''}

\begin{ucspec}{Đặc tả use case ``Quản lý mã giảm giá''}
\textbf{Mã use case} & UC27 \\ \hline
\textbf{Tên use case} & Quản lý mã giảm giá \\ \hline
\textbf{Tác nhân} & Quản trị viên \\ \hline
\textbf{Mục đích} & Tạo và quản lý các mã giảm giá (voucher) \\ \hline
\textbf{Tiền điều kiện} & Quản trị viên đã đăng nhập \\ \hline
\textbf{Hậu điều kiện} & Voucher được tạo/cập nhật \\ \hline
\textbf{Luồng sự kiện chính} &
1. Quản trị viên vào trang quản lý mã giảm giá. \newline
2. Tạo mã mới: nhập mã, loại giảm (\% hoặc số tiền), giá trị, điều kiện đơn tối thiểu, giới hạn
sử dụng, thời gian hiệu lực. \newline
3. Hệ thống kiểm tra trùng mã và lưu voucher. \\ \hline
\textbf{Luồng thay thế} &
3a. Mã đã tồn tại: hệ thống thông báo lỗi. \\ \hline
\end{ucspec}

\subsection{Đặc tả use case ``Kiểm duyệt đánh giá''}

\begin{ucspec}{Đặc tả use case ``Kiểm duyệt đánh giá''}
\textbf{Mã use case} & UC31 \\ \hline
\textbf{Tên use case} & Kiểm duyệt đánh giá \\ \hline
\textbf{Tác nhân} & Quản trị viên \\ \hline
\textbf{Mục đích} & Duyệt, từ chối hoặc phản hồi các đánh giá của khách hàng \\ \hline
\textbf{Tiền điều kiện} & Có đánh giá đang ở trạng thái chờ duyệt \\ \hline
\textbf{Hậu điều kiện} & Trạng thái đánh giá được cập nhật; điểm trung bình sản phẩm được tính lại \\ \hline
\textbf{Luồng sự kiện chính} &
1. Quản trị viên vào trang kiểm duyệt đánh giá. \newline
2. Xem nội dung, hình ảnh đánh giá. \newline
3. Chọn duyệt / từ chối / phản hồi / xóa. \newline
4. Hệ thống cập nhật trạng thái và tính lại điểm trung bình của sản phẩm. \\ \hline
\textbf{Luồng thay thế} & --- \\ \hline
\end{ucspec}

\subsection{Đặc tả use case ``Thống kê báo cáo doanh thu''}

\begin{ucspec}{Đặc tả use case ``Thống kê báo cáo doanh thu''}
\textbf{Mã use case} & UC32 \\ \hline
\textbf{Tên use case} & Thống kê, báo cáo doanh thu \\ \hline
\textbf{Tác nhân} & Quản trị viên \\ \hline
\textbf{Mục đích} & Theo dõi doanh thu và các chỉ số kinh doanh \\ \hline
\textbf{Tiền điều kiện} & Quản trị viên đã đăng nhập \\ \hline
\textbf{Hậu điều kiện} & Hiển thị báo cáo; có thể xuất tệp Excel \\ \hline
\textbf{Luồng sự kiện chính} &
1. Quản trị viên vào trang báo cáo. \newline
2. Hệ thống tổng hợp doanh thu theo ngày, sản phẩm bán chạy, khách hàng và các chỉ số khác. \newline
3. Hệ thống hiển thị biểu đồ và bảng số liệu. \newline
4. Quản trị viên xuất báo cáo ra tệp Excel khi cần. \\ \hline
\textbf{Luồng thay thế} & --- \\ \hline
\end{ucspec}

Dữ liệu đầu vào của use case ``Quản lý mã giảm giá'' được mô tả trong Bảng~\ref{tab:voucher-input}
(theo lược đồ kiểm tra dữ liệu thực tế của hệ thống).

\begin{longtable}{|c|>{\raggedright\arraybackslash}p{0.24\linewidth}|c|>{\raggedright\arraybackslash}p{0.34\linewidth}|}
\caption{Dữ liệu đầu vào use case ``Quản lý mã giảm giá''}\label{tab:voucher-input}\\
\hline
\rowcolor{gray!12}\textbf{STT} & \textbf{Trường} & \textbf{Bắt buộc} & \textbf{Điều kiện hợp lệ} \\ \hline
\endfirsthead
\hline
\rowcolor{gray!12}\textbf{STT} & \textbf{Trường} & \textbf{Bắt buộc} & \textbf{Điều kiện hợp lệ} \\ \hline
\endhead
1 & Mã giảm giá & Có & 3--30 ký tự, chỉ chữ HOA, số, gạch ngang/dưới \\ \hline
2 & Mô tả & Không & Tối đa 200 ký tự \\ \hline
3 & Loại giảm & Có & PERCENT hoặc FIXED \\ \hline
4 & Giá trị giảm & Có & Số nguyên dương; nếu \% phải từ 1--100 \\ \hline
5 & Đơn tối thiểu & Không & Số nguyên không âm \\ \hline
6 & Giảm tối đa & Không & Số nguyên dương \\ \hline
7 & Giới hạn lượt dùng & Không & Số nguyên dương \\ \hline
8 & Giới hạn theo người dùng & Không & Số nguyên dương \\ \hline
9 & Thời gian hiệu lực & Có & Ngày kết thúc phải sau ngày bắt đầu \\ \hline
\end{longtable}

\subsection{Đặc tả use case ``Đăng nhập quản trị''}

\begin{ucspec}{Đặc tả use case ``Đăng nhập quản trị''}
\textbf{Mã use case} & UC21 \\ \hline
\textbf{Tên use case} & Đăng nhập quản trị \\ \hline
\textbf{Tác nhân} & Quản trị viên \\ \hline
\textbf{Mục đích} & Truy cập khu vực quản trị với quyền quản trị \\ \hline
\textbf{Tiền điều kiện} & Tài khoản có vai trò ADMIN \\ \hline
\textbf{Hậu điều kiện} & Quản trị viên truy cập được trang quản trị \\ \hline
\textbf{Luồng sự kiện chính} &
1. Quản trị viên đăng nhập bằng email và mật khẩu. \newline
2. Hệ thống xác thực và kiểm tra vai trò người dùng. \newline
3. Nếu là ADMIN, cho phép truy cập khu vực quản trị. \\ \hline
\textbf{Luồng thay thế} &
3a. Tài khoản không phải ADMIN: từ chối truy cập khu vực quản trị. \\ \hline
\end{ucspec}

\subsection{Đặc tả use case ``Quản lý danh mục sản phẩm''}

\begin{ucspec}{Đặc tả use case ``Quản lý danh mục sản phẩm''}
\textbf{Mã use case} & UC22 \\ \hline
\textbf{Tên use case} & Quản lý danh mục sản phẩm \\ \hline
\textbf{Tác nhân} & Quản trị viên \\ \hline
\textbf{Mục đích} & Thêm, sửa, xóa danh mục sản phẩm (có phân cấp cha--con) \\ \hline
\textbf{Tiền điều kiện} & Quản trị viên đã đăng nhập \\ \hline
\textbf{Hậu điều kiện} & Cây danh mục được cập nhật \\ \hline
\textbf{Luồng sự kiện chính} &
1. Quản trị viên vào trang quản lý danh mục. \newline
2. Thêm/sửa danh mục: nhập tên, đường dẫn (slug), chọn danh mục cha, thứ tự hiển thị. \newline
3. Hệ thống kiểm tra trùng slug và lưu danh mục. \\ \hline
\textbf{Luồng thay thế} &
3a. Slug đã tồn tại: hệ thống thông báo lỗi. \\ \hline
\end{ucspec}

\subsection{Đặc tả use case ``Quản lý tồn kho''}

\begin{ucspec}{Đặc tả use case ``Quản lý tồn kho''}
\textbf{Mã use case} & UC24 \\ \hline
\textbf{Tên use case} & Quản lý tồn kho \\ \hline
\textbf{Tác nhân} & Quản trị viên \\ \hline
\textbf{Mục đích} & Điều chỉnh và theo dõi tồn kho của các biến thể \\ \hline
\textbf{Tiền điều kiện} & Quản trị viên đã đăng nhập \\ \hline
\textbf{Hậu điều kiện} & Tồn kho được cập nhật; thay đổi được ghi nhật ký \\ \hline
\textbf{Luồng sự kiện chính} &
1. Quản trị viên vào trang quản lý tồn kho. \newline
2. Điều chỉnh số lượng tồn của một biến thể (nhập thêm/bớt) kèm lý do. \newline
3. Hệ thống cập nhật tồn kho và ghi một bản ghi vào nhật ký tồn kho. \\ \hline
\textbf{Luồng thay thế} & --- \\ \hline
\end{ucspec}

\subsection{Đặc tả use case ``Xử lý đổi/trả hàng''}

\begin{ucspec}{Đặc tả use case ``Xử lý đổi/trả hàng''}
\textbf{Mã use case} & UC26 \\ \hline
\textbf{Tên use case} & Xử lý đổi/trả hàng \\ \hline
\textbf{Tác nhân} & Quản trị viên \\ \hline
\textbf{Mục đích} & Duyệt, từ chối hoặc hoàn tiền cho yêu cầu đổi/trả \\ \hline
\textbf{Tiền điều kiện} & Có yêu cầu đổi/trả đang chờ xử lý \\ \hline
\textbf{Hậu điều kiện} & Trạng thái yêu cầu được cập nhật \\ \hline
\textbf{Luồng sự kiện chính} &
1. Quản trị viên xem danh sách yêu cầu đổi/trả. \newline
2. Chọn một yêu cầu để xem chi tiết. \newline
3. Duyệt, từ chối hoặc xác nhận hoàn tiền. \newline
4. Khi hoàn tiền: hệ thống hoàn trả tồn kho và đánh dấu đơn ``đã hoàn tiền''. \\ \hline
\textbf{Luồng thay thế} &
3a. Từ chối: cập nhật trạng thái từ chối kèm ghi chú. \\ \hline
\end{ucspec}

\subsection{Đặc tả use case ``Quản lý flash sale''}

\begin{ucspec}{Đặc tả use case ``Quản lý flash sale''}
\textbf{Mã use case} & UC28 \\ \hline
\textbf{Tên use case} & Quản lý flash sale \\ \hline
\textbf{Tác nhân} & Quản trị viên \\ \hline
\textbf{Mục đích} & Tạo và quản lý chương trình giảm giá theo thời gian \\ \hline
\textbf{Tiền điều kiện} & Quản trị viên đã đăng nhập \\ \hline
\textbf{Hậu điều kiện} & Chương trình flash sale được tạo/cập nhật \\ \hline
\textbf{Luồng sự kiện chính} &
1. Quản trị viên vào trang quản lý flash sale. \newline
2. Tạo chương trình: đặt tên, thời gian bắt đầu/kết thúc, chọn sản phẩm và giá khuyến mãi. \newline
3. Hệ thống lưu chương trình; giá khuyến mãi tự áp dụng khi chương trình đang hiệu lực. \\ \hline
\textbf{Luồng thay thế} & --- \\ \hline
\end{ucspec}

\subsection{Đặc tả use case ``Quản lý banner và nội dung''}

\begin{ucspec}{Đặc tả use case ``Quản lý banner và nội dung''}
\textbf{Mã use case} & UC29 \\ \hline
\textbf{Tên use case} & Quản lý banner và nội dung \\ \hline
\textbf{Tác nhân} & Quản trị viên \\ \hline
\textbf{Mục đích} & Quản lý các banner quảng cáo hiển thị trên trang chủ \\ \hline
\textbf{Tiền điều kiện} & Quản trị viên đã đăng nhập \\ \hline
\textbf{Hậu điều kiện} & Danh sách và thứ tự banner được cập nhật \\ \hline
\textbf{Luồng sự kiện chính} &
1. Quản trị viên vào trang quản lý banner. \newline
2. Thêm/sửa banner: nhập tiêu đề, ảnh, liên kết, trạng thái hiển thị. \newline
3. Sắp xếp thứ tự banner bằng thao tác kéo thả. \newline
4. Hệ thống lưu và phản ánh ngay trên trang chủ. \\ \hline
\textbf{Luồng thay thế} & --- \\ \hline
\end{ucspec}

\subsection{Đặc tả use case ``Quản lý tài khoản khách hàng''}

\begin{ucspec}{Đặc tả use case ``Quản lý tài khoản khách hàng''}
\textbf{Mã use case} & UC30 \\ \hline
\textbf{Tên use case} & Quản lý tài khoản khách hàng \\ \hline
\textbf{Tác nhân} & Quản trị viên \\ \hline
\textbf{Mục đích} & Xem danh sách và thông tin chi tiết khách hàng \\ \hline
\textbf{Tiền điều kiện} & Quản trị viên đã đăng nhập \\ \hline
\textbf{Hậu điều kiện} & Hiển thị thông tin khách hàng \\ \hline
\textbf{Luồng sự kiện chính} &
1. Quản trị viên vào trang quản lý khách hàng. \newline
2. Tìm kiếm và xem danh sách khách hàng. \newline
3. Chọn một khách hàng để xem chi tiết (thông tin, lịch sử đơn hàng). \\ \hline
\textbf{Luồng thay thế} & --- \\ \hline
\end{ucspec}

\section{Yêu cầu phi chức năng}

\subsection{Yêu cầu về bảo mật}
\begin{itemize}
  \item Mật khẩu người dùng được quản lý và mã hóa bởi dịch vụ xác thực Supabase Auth.
  \item Mọi dữ liệu đầu vào được kiểm tra và làm sạch ở phía máy chủ bằng lược đồ Zod.
  \item Phân quyền truy cập: chỉ quản trị viên mới truy cập được khu vực quản trị; áp dụng thêm
  Row-Level Security ở tầng cơ sở dữ liệu.
  \item Webhook thanh toán được xác thực bằng chữ ký HMAC; giới hạn tần suất truy cập (rate
  limit) cho các thao tác nhạy cảm.
\end{itemize}

\subsection{Yêu cầu về hiệu năng}
\begin{itemize}
  \item Trang sản phẩm và trang chủ được tối ưu bằng React Server Components và bộ nhớ đệm
  (caching) để có thời gian phản hồi nhanh.
  \item Truy vấn cơ sở dữ liệu được tối ưu bằng chỉ mục (index) trên các cột truy vấn thường xuyên.
\end{itemize}

\subsection{Yêu cầu về khả năng bảo trì}
\begin{itemize}
  \item Mã nguồn viết bằng TypeScript, tổ chức theo tầng (giao diện, nghiệp vụ, truy vấn dữ liệu).
  \item Logic kiểm tra dữ liệu được tập trung và tái sử dụng; có bộ kiểm thử tự động.
\end{itemize}

\subsection{Yêu cầu về khả năng mở rộng}
\begin{itemize}
  \item Kiến trúc cho phép bổ sung danh mục, sản phẩm và tính năng mới mà không ảnh hưởng lớn
  tới hệ thống hiện có.
  \item Triển khai serverless trên Vercel cho phép co giãn theo lưu lượng.
\end{itemize}

\subsection{Yêu cầu về tính khả dụng và giao diện}
\begin{itemize}
  \item Giao diện thân thiện, hiển thị tốt trên cả máy tính và thiết bị di động (responsive).
  \item Thông báo lỗi rõ ràng, dẫn dắt người dùng hoàn thành thao tác.
\end{itemize}

\subsection{Yêu cầu về cơ sở dữ liệu}
\begin{itemize}
  \item Dữ liệu được lưu trữ nhất quán trên PostgreSQL; các nghiệp vụ nhạy cảm (đặt hàng, trừ
  kho, áp dụng voucher) được thực hiện trong giao dịch để bảo đảm toàn vẹn.
  \item Giá trị tiền tệ được lưu dưới dạng số nguyên (đồng) để tránh sai số dấu phẩy động.
\end{itemize}

\subsection{Ma trận phân quyền}

Bảng~\ref{tab:rbac} tóm tắt quyền truy cập của từng nhóm chức năng theo vai trò người dùng.

\begin{longtable}{|>{\raggedright\arraybackslash}p{0.46\linewidth}|c|c|}
\caption{Ma trận phân quyền người dùng}\label{tab:rbac}\\
\hline
\rowcolor{gray!12}\textbf{Chức năng} & \textbf{Khách hàng} & \textbf{Quản trị viên} \\ \hline
\endfirsthead
\hline
\rowcolor{gray!12}\textbf{Chức năng} & \textbf{Khách hàng} & \textbf{Quản trị viên} \\ \hline
\endhead
Duyệt, tìm kiếm, xem sản phẩm & \checkmark & \checkmark \\ \hline
Giỏ hàng, yêu thích, so sánh & \checkmark & \\ \hline
Đặt hàng, thanh toán, đổi/trả & \checkmark & \\ \hline
Đánh giá sản phẩm & \checkmark & \\ \hline
Quản lý sản phẩm, danh mục, kho & & \checkmark \\ \hline
Quản lý đơn hàng & & \checkmark \\ \hline
Quản lý khuyến mãi, banner & & \checkmark \\ \hline
Kiểm duyệt đánh giá & & \checkmark \\ \hline
Quản lý khách hàng & & \checkmark \\ \hline
Thống kê, báo cáo & & \checkmark \\ \hline
\end{longtable}
